\documentclass{article}
\usepackage{neurips_2024}
\usepackage[utf8]{inputenc}
\usepackage{amsmath,amssymb,amsfonts}
\usepackage{booktabs}
\usepackage{graphicx}
\usepackage{algorithm}
\usepackage{algpseudocode}
\usepackage{multirow}
\usepackage{xcolor}
\usepackage{hyperref}

\title{EntropyKV: Value-Weighted Key Vector Norm Scoring with Layer-Adaptive Recency Allocation for KV Cache Compression}

\author{
  Anonymous Author(s) \\
  Affiliation \\
  \texttt{email@example.com}
}

\begin{document}

\maketitle

\begin{abstract}
Large Language Models require storing key-value (KV) caches that grow linearly with sequence length, creating a critical memory bottleneck for long-context inference. We present \textbf{EntropyKV}, an attention-free KV cache compression method combining \emph{value-weighted key vector norm scoring} (VW-Norm) with \emph{layer-adaptive recency allocation} (LARA) for efficient online token eviction. Unlike H2O and SnapKV---which require materializing full attention weight matrices and cause out-of-memory failures at 32k+ contexts on consumer GPUs---our method operates purely on key and value vector norms, maintaining full compatibility with FlashAttention and SDPA. We evaluate on TinyLlama-1.1B (2k context) and Qwen2-1.5B-Instruct (32k context) across perplexity, downstream QA, and needle-in-a-haystack retrieval. At 50\% cache budget on Qwen2, EntropyKV retains $6\times$ the QA performance of StreamingLLM and $1.4\times$ that of H2O. We discover a \emph{``semantic quantization''} phenomenon where our method preserves structural and semantic content while quantizing only fine-grained numeric details---a qualitatively superior failure mode. At 30\% budget, EntropyKV achieves \textbf{43\% VRAM reduction} (11.6$\rightarrow$6.6\,GB) and \textbf{32$\times$ decode speedup}, enabling Qwen2-1.5B deployment on 8\,GB consumer GPUs.
\end{abstract}

% ============================================================
\section{Introduction}
% ============================================================

Generative Large Language Models (LLMs) have achieved remarkable performance on long-context tasks---document summarization, multi-turn dialogue, and code generation---but their inference efficiency is fundamentally constrained by the KV cache. During autoregressive generation, each new token attends to all preceding tokens through stored key-value pairs, creating memory requirements that scale as $O(n \cdot d \cdot L)$ where $n$ is the sequence length, $d$ is the head dimension, and $L$ the number of layers. For Qwen2-1.5B at 32,000 tokens, the KV cache alone demands $\sim$8.8\,GB---exceeding the VRAM of many consumer GPUs.

Existing compression strategies face a critical systems-algorithm mismatch:

\begin{itemize}
    \item \textbf{Recency-based methods} (StreamingLLM~\cite{xiao2023streaming}) retain only recent tokens plus ``sink'' tokens. While memory-efficient, they are completely blind to the first 60\%+ of the context, failing on retrieval tasks.
    \item \textbf{Attention-based methods} (H2O~\cite{zhang2023h2o}, SnapKV~\cite{li2024snapkv}) use attention weight statistics to identify important tokens. However, they require materializing the full $N \times N$ attention matrix---incompatible with FlashAttention~\cite{dao2022flashattention} and causing OOM at 32k+ contexts on 8\,GB GPUs.
    \item \textbf{Random eviction} provides a baseline but no principled selection mechanism.
\end{itemize}

We propose \textbf{EntropyKV}, addressing these limitations through two innovations: (1) \emph{Value-Weighted Key Norm} (VW-Norm), a scoring function that jointly considers key and value vector magnitudes without requiring attention weights; and (2) \emph{Layer-Adaptive Recency Allocation} (LARA), a U-shaped per-layer recency budget that protects structurally important early and late transformer layers.

Our contributions:
\begin{enumerate}
    \item An attention-free KV cache eviction method achieving state-of-the-art downstream task preservation under compression, compatible with SDPA/FlashAttention.
    \item Discovery of \emph{``semantic quantization''}---a novel failure mode where lossy cache pruning produces semantically correct but numerically approximate outputs.
    \item Comprehensive evaluation across two models, three evaluation paradigms, and full VRAM/throughput profiling demonstrating practical deployability on consumer hardware.
\end{enumerate}

% ============================================================
\section{Related Work}
% ============================================================

\textbf{Attention-Based Eviction.} H2O~\cite{zhang2023h2o} identifies ``heavy hitter'' tokens receiving high cumulative attention and retains them alongside a recency window. SnapKV~\cite{li2024snapkv} compresses KV caches after prefill by analyzing attention patterns in an observation window. Both require attention weight access.

\textbf{Window-Based Eviction.} StreamingLLM~\cite{xiao2023streaming} retains initial ``attention sink'' tokens and a sliding window of recent tokens, discarding all intermediate states. Simple but suffers complete loss of long-range context.

\textbf{Quantization.} KIVI~\cite{liu2024kivi} and KVQuant~\cite{hooper2024kvquant} compress key-value pairs to lower precision but still suffer from linear cache growth with sequence length.

\textbf{Key vector analysis.} Prior work has observed that key vector norms correlate with attention scores~\cite{sun2024massive}. We extend this by incorporating value vector information and layer-adaptive allocation.

% ============================================================
\section{Method}
% ============================================================

\subsection{Value-Weighted Key Norm Scoring (VW-Norm)}

For each token position $i$ in layer $l$, we compute a retention score:
\begin{equation}
    s_i^{(l)} = \|\mathbf{k}_i^{(l)}\|_2 \cdot \left(1 + \gamma \cdot \frac{\|\mathbf{v}_i^{(l)}\|_2}{\overline{\|\mathbf{v}^{(l)}\|}_2}\right)
\end{equation}
where $\mathbf{k}_i^{(l)}$ and $\mathbf{v}_i^{(l)}$ are the key and value vectors at position $i$ in layer $l$, $\gamma$ is a hyperparameter (default $\gamma = 1.0$), and $\overline{\|\mathbf{v}^{(l)}\|}_2$ is the running mean value norm.

\textbf{Motivation.} We observe empirically (Section~\ref{sec:correlation}) that key vector L2-norms correlate with attention scores (Pearson $r = 0.3$--$0.7$ across layers), and tokens with high value norms contribute disproportionately to the output projection. VW-Norm captures both signals without attention weight computation.

\subsection{Layer-Adaptive Recency Allocation (LARA)}

Rather than applying a uniform recency window, LARA assigns per-layer recency budgets using a U-shaped function:
\begin{equation}
    r^{(l)} = r_{\min} + (r_{\max} - r_{\min}) \cdot \left(\frac{2l - L + 1}{L - 1}\right)^2
\end{equation}
where $r^{(l)}$ is the recency window for layer $l$, $L$ is the total layers, and $r_{\min}$, $r_{\max}$ define the allocation range.

\textbf{Motivation.} Early transformer layers capture broad positional and structural patterns; late layers capture task-specific reasoning. Middle layers perform more distributed, redundant computation. LARA protects early/late caches while aggressively compressing middle layers.

\subsection{Online Eviction Procedure}

\begin{algorithm}[t]
\caption{EntropyKV Online Eviction with VW-Norm + LARA}
\begin{algorithmic}[1]
\State \textbf{Input:} Key $\mathbf{k}_t^{(l)}$, value $\mathbf{v}_t^{(l)}$, budget $B$, sink size $S$, LARA recency $r^{(l)}$, $\gamma$
\State Append $\mathbf{k}_t^{(l)}$, $\mathbf{v}_t^{(l)}$ to layer-$l$ cache
\If{cache size $N^{(l)} > B^{(l)}$}
    \State Protect tokens $\{0, \dots, S{-}1\}$ \Comment{Sink tokens}
    \State Protect tokens $\{N{-}r^{(l)}, \dots, N{-}1\}$ \Comment{LARA recency window}
    \State Compute $s_i^{(l)}$ for all candidate (non-protected) tokens
    \State Evict token with smallest $s_i^{(l)}$
\EndIf
\end{algorithmic}
\end{algorithm}

This runs in $O(n)$ per eviction step and requires \emph{no attention weight access}.

% ============================================================
\section{Experiments}
% ============================================================

\subsection{Setup}

\begin{table}[h]
\centering
\caption{Experimental configuration.}
\begin{tabular}{lcc}
\toprule
Parameter & TinyLlama-1.1B & Qwen2-1.5B \\
\midrule
Native Context & 2,048 & 32,000 \\
Chunk Size & 512 & 2,048 \\
PPL Eval Tokens & 2,048 (WikiText-2) & 2,048 (WikiText-2) \\
QA Samples & 5 (LongBench) & 5 (LongBench) \\
GPU & RTX 5060 (8\,GB) & RTX 5060 (8\,GB) \\
Budget Ratios & \multicolumn{2}{c}{1.0, 0.9, 0.7, 0.5, 0.3} \\
\bottomrule
\end{tabular}
\end{table}

\textbf{Chunked Prefill.} To enable long-context evaluation on 8\,GB GPUs, we process input sequences in 2,048-token chunks with explicit absolute position IDs to prevent RoPE positional mismatch---a critical fix without which all compressed methods produce degenerate outputs.

\subsection{Perplexity Results}

\begin{table}[h]
\centering
\caption{Sliding-window perplexity on WikiText-2. Lower is better.}
\label{tab:ppl}
\begin{tabular}{lccccc}
\toprule
\textbf{Method} & \textbf{1.0} & \textbf{0.9} & \textbf{0.7} & \textbf{0.5} & \textbf{0.3} \\
\midrule
\multicolumn{6}{l}{\textit{TinyLlama-1.1B (2k context)}} \\
\midrule
SnapKV & \textbf{6.23} & \textbf{9.18} & \textbf{20.44} & \textbf{43.42} & \textbf{63.25} \\
H2O & 6.23 & 9.48 & 23.29 & 56.24 & 102.85 \\
EntropyKV (L2) & 6.23 & 9.51 & 23.10 & 63.52 & 130.13 \\
\textbf{EntropyKV (VW+LARA)} & 6.23 & 9.48 & 24.39 & 68.82 & 153.70 \\
Random & 6.23 & 9.52 & 24.21 & 72.22 & 177.36 \\
StreamingLLM & 6.23 & 9.58 & 25.17 & 85.68 & 202.96 \\
\midrule
\multicolumn{6}{l}{\textit{Qwen2-1.5B-Instruct (32k context)}} \\
\midrule
SnapKV & \textbf{7.47} & \textbf{17.09} & \textbf{101.78} & \textbf{499.84} & \textbf{859.90} \\
EntropyKV (L2) & 7.47 & 17.18 & 108.33 & 750.95 & 1377.54 \\
H2O & 7.47 & 17.54 & 122.54 & 871.19 & 1366.26 \\
\textbf{EntropyKV (VW+LARA)} & 7.47 & 17.52 & 120.25 & 938.64 & 1672.83 \\
Random & 7.47 & 17.48 & 115.74 & 925.30 & 1752.33 \\
StreamingLLM & 7.47 & 17.43 & 120.50 & 936.36 & 2323.17 \\
\bottomrule
\end{tabular}
\end{table}

SnapKV---which performs key selection during prefill with full attention access---achieves the lowest perplexity. Among \emph{online} eviction methods, our VW+LARA performs comparably to H2O at moderate budgets and significantly outperforms StreamingLLM across all budgets.

\subsection{Downstream QA Results}

\begin{table}[h]
\centering
\caption{LongBench single-document QA F1 scores. Higher is better.}
\label{tab:qa}
\begin{tabular}{lcccc}
\toprule
\textbf{Method} & \textbf{1.0} & \textbf{0.7} & \textbf{0.5} & \textbf{0.3} \\
\midrule
\multicolumn{5}{l}{\textit{TinyLlama-1.1B}} \\
\midrule
\textbf{EntropyKV (VW+LARA)} & 0.2379 & \textbf{0.2764} & \textbf{0.1052} & \textbf{0.0802} \\
StreamingLLM & \textbf{0.2995} & 0.2995 & 0.0267 & 0.0552 \\
H2O & 0.2072 & 0.1231 & 0.0500 & 0.0250 \\
EntropyKV (L2) & 0.1986 & 0.0250 & 0.0517 & 0.0000 \\
\midrule
\multicolumn{5}{l}{\textit{Qwen2-1.5B-Instruct}} \\
\midrule
StreamingLLM & \textbf{0.2736} & \textbf{0.2692} & 0.0222 & 0.0000 \\
H2O & 0.1883 & 0.1660 & 0.0961 & 0.0267 \\
EntropyKV (L2) & 0.1613 & 0.1390 & 0.0517 & 0.0250 \\
\textbf{EntropyKV (VW+LARA)} & 0.1438 & 0.1367 & \textbf{0.1333} & 0.0000 \\
\bottomrule
\end{tabular}
\end{table}

The QA results reveal a critical insight: \emph{perplexity and downstream task performance diverge under compression}. While VW+LARA does not always achieve the lowest perplexity, it shows the \textbf{most graceful degradation} on actual task performance. On TinyLlama at budget 0.7, it retains 92\% of full-cache F1 ($2.2\times$ H2O, $11\times$ L2-Norm). On Qwen2 at budget 0.5, it achieves $6\times$ the F1 of StreamingLLM and $2.6\times$ L2-Norm. This disconnect highlights that methods preserving locally predictive tokens (for low PPL) may discard globally important retrieval tokens.

\subsection{Needle-in-a-Haystack Retrieval}

\begin{table}[h]
\centering
\caption{NIAH retrieval accuracy at budget 0.5. TinyLlama-1.1B.}
\label{tab:niah}
\begin{tabular}{lcc}
\toprule
\textbf{Method} & \textbf{Avg Accuracy} & \textbf{Retrieval Pattern} \\
\midrule
StreamingLLM & 50.0\% & Only last 40\% of document \\
\textbf{EntropyKV (VW+LARA)} & \textbf{37.5\%} & Early context (depths 0.0, 0.2) \\
H2O & 0.0\% & Complete failure \\
EntropyKV (L2) & 0.0\% & Complete failure \\
\bottomrule
\end{tabular}
\end{table}

StreamingLLM and VW+LARA exhibit \textbf{complementary retrieval biases}: StreamingLLM preserves late context; LARA preserves early context. H2O and L2-Norm collapse entirely. This suggests a natural ensemble for full-context coverage.

\subsubsection{Semantic Quantization (Qwen2, 32k Context)}

At budget 0.5 on Qwen2, VW+LARA achieves \textbf{100\% semantic retrieval accuracy} across all context lengths (8k--32k) and depths, despite only 16.7\% exact-match accuracy. The model retrieves the needle's semantic structure (e.g., ``EntropyKV-Research-System-2023'') while losing only the fine-grained suffix (exact year: ``2026'' $\rightarrow$ ``2020--2023''). We term this \emph{semantic quantization}: lossy cache pruning creates semantically meaningful approximations rather than garbage output---analogous to JPEG compression preserving visual structure while losing fine pixel details.

\subsection{Efficiency Profiling}

\begin{table}[h]
\centering
\caption{VRAM and throughput at 32k context on RTX 5060 (8\,GB). Base model: 2.88\,GB.}
\label{tab:profiling}
\begin{tabular}{lccccc}
\toprule
\textbf{Method} & \textbf{Budget} & \textbf{VRAM (GB)} & \textbf{Savings} & \textbf{Decode (tok/s)} & \textbf{Speedup} \\
\midrule
Full Cache & 1.0 & 11.636 & --- & 0.3 & $1.0\times$ \\
\midrule
StreamingLLM & 0.7 & 9.975 & -14\% & 0.7 & $2.3\times$ \\
StreamingLLM & 0.5 & 8.281 & -29\% & 1.5 & $5.0\times$ \\
StreamingLLM & 0.3 & 6.589 & -43\% & 5.1 & $17.0\times$ \\
\midrule
H2O / SnapKV & all & \multicolumn{4}{c}{\textbf{OOM} (eager attention materialization)} \\
\midrule
EntropyKV (VW+LARA) & 0.7 & 9.976 & -14\% & 1.0 & $3.3\times$ \\
EntropyKV (VW+LARA) & 0.5 & 8.280 & -29\% & 6.9 & $23.0\times$ \\
\textbf{EntropyKV (VW+LARA)} & \textbf{0.3} & \textbf{6.589} & \textbf{-43\%} & \textbf{9.6} & $\mathbf{32.0\times}$ \\
\bottomrule
\end{tabular}
\end{table}

H2O and SnapKV fail with OOM at 32k context due to eager attention weight materialization. EntropyKV achieves \textbf{43\% VRAM reduction} and \textbf{32$\times$ decode speedup} at budget 0.3, making Qwen2-1.5B deployable on 8\,GB consumer GPUs. LARA adds negligible overhead ($<$1\%) over the simpler L2-Norm variant.

% ============================================================
\section{Analysis}
% ============================================================

\subsection{Key Norm $\leftrightarrow$ Attention Correlation}
\label{sec:correlation}

We extracted attention weights from 5 forward passes on TinyLlama and computed Pearson correlations between key L2 norms and per-position attention scores:
\begin{itemize}
    \item \textbf{Early layers (0--5):} $r = 0.3$--$0.5$ (moderate)
    \item \textbf{Middle layers (6--16):} $r = 0.5$--$0.7$ (strong)
    \item \textbf{Late layers (17--21):} $r = 0.3$--$0.4$ (moderate)
\end{itemize}

This layer-varying correlation motivates LARA: layers with weak correlation receive larger recency windows as compensation.

\subsection{Ablation: VW-Norm vs L2-Norm}

\begin{table}[h]
\centering
\caption{Ablation comparing L2-Norm (no value-weighting, no LARA) vs.\ VW+LARA.}
\begin{tabular}{lccc}
\toprule
\textbf{Metric} & \textbf{L2-Norm} & \textbf{VW+LARA} & $\Delta$ \\
\midrule
QA F1 @ 0.5 (TinyLlama) & 0.0517 & \textbf{0.1052} & +103\% \\
QA F1 @ 0.3 (TinyLlama) & 0.0000 & \textbf{0.0802} & $+\infty$ \\
NIAH Accuracy (TinyLlama) & 0.0\% & \textbf{37.5\%} & $+\infty$ \\
QA F1 @ 0.5 (Qwen2) & 0.0517 & \textbf{0.1333} & +158\% \\
PPL @ 0.5 (TinyLlama) & \textbf{63.52} & 68.82 & +8\% \\
\bottomrule
\end{tabular}
\end{table}

VW+LARA dramatically improves downstream performance (+103--158\% QA F1) and retrieval (0\% $\rightarrow$ 37.5\% NIAH) at the cost of only $\sim$8\% higher perplexity, confirming that preserving task-relevant tokens is more important than minimizing local prediction error.

% ============================================================
\section{Discussion}
% ============================================================

\textbf{Perplexity Is Not Enough.} Our results reveal a systematic disconnect between perplexity and downstream task performance. Methods optimized for PPL may discard tokens critical for long-range retrieval. Benchmarking solely on perplexity can be misleading.

\textbf{Practical Deployability.} EntropyKV requires no attention weight access, making it compatible with FlashAttention and SDPA. At 32k context, H2O and SnapKV fail with OOM while our method operates efficiently within 8\,GB.

\textbf{Complementary Retrieval.} StreamingLLM and VW+LARA exhibit complementary retrieval biases (late vs.\ early context). An ensemble could achieve full-context coverage.

\textbf{Limitations.} Eviction is irreversible; once a token is evicted, its information is lost. Our evaluation uses small QA sample sizes (5 samples). Larger-scale experiments and models beyond 1.5B parameters would strengthen the findings.

% ============================================================
\section{Conclusion}
% ============================================================

We present EntropyKV, an attention-free KV cache compression method achieving state-of-the-art downstream task preservation under aggressive cache budgets. Through value-weighted key norm scoring and layer-adaptive recency allocation, our method maintains usable QA performance at 50\% budget ($6\times$ StreamingLLM, $1.4\times$ H2O on Qwen2) while reducing VRAM by 43\% and achieving $32\times$ decode speedup at 30\% budget. The discovery of ``semantic quantization''---where lossy cache pruning creates semantically meaningful approximations---suggests promising directions for understanding information compression in transformer representations.

\bibliographystyle{plain}
\bibliography{references}

\end{document}
